\section{Motivation}

Women in dual-earner couples work systematically fewer hours than their male partners and earn lower wages, even conditional on education \citep{blau2017gender}.
\citet{kleven2019children} show that childbirth explains a large share of this gap, but a substantial residual persists among couples without children.
A structural explanation is the tax system: under joint filing the secondary earner faces an effective marginal rate equal to the household's top bracket, which is typically determined by the primary earner's (higher) income.
This concentrates the labor supply distortion on the partner with lower initial earnings, reinforcing specialization over the life cycle through reduced human capital accumulation.

Several studies document large labor supply elasticities for secondary earners \citep{bick2017taxation}, making the tax wedge a plausible driver of the gender gap.
Yet most analyses are static and miss the dynamic human capital channel: lower hours today reduce future wages, which in turn suppress future labor supply --- a self-reinforcing mechanism that a dynamic model can quantify.
